\documentclass[12pt,a4paper]{article}
\usepackage[utf8]{inputenc}
\usepackage[vietnamese]{babel}
\usepackage{amsmath,amssymb,amsthm}
\usepackage{geometry}
\usepackage{enumitem}
\usepackage[most]{tcolorbox}
\usepackage{hyperref}
\usepackage{fancyhdr}
\usepackage{tikz}
\usepackage{eso-pic}
\usepackage{xcolor}
\geometry{a4paper, margin=2cm, bottom=2.5cm}
\hypersetup{colorlinks=true, linkcolor=blue!70!black}
\setlist[itemize]{topsep=4pt, itemsep=2pt}
\setlist[enumerate]{topsep=4pt, itemsep=2pt}
\AddToShipoutPictureFG{%
  \AtPageCenter{%
    \begin{tikzpicture}[remember picture, overlay]
      \node[rotate=45, scale=1, opacity=0.06]{\fontsize{60}{72}\selectfont\bfseries\sffamily Đặng Tấn Quí};
    \end{tikzpicture}%
  }%
}
\pagestyle{fancy}
\fancyhf{}
\fancyhead[L]{\small\textit{LÝ THUYẾT ĐỘ ĐO VÀ TÍCH PHÂN}}
\fancyhead[R]{\small\textit{Đặng Tấn Quí}}
\fancyfoot[C]{\thepage}
\fancyfoot[L]{\footnotesize\textcopyright\ Đặng Tấn Quí}
\fancyfoot[R]{\footnotesize SĐT: 0377\,122\,210}
\renewcommand{\headrulewidth}{0.4pt}
\renewcommand{\footrulewidth}{0.4pt}
\fancypagestyle{plain}{\fancyhf{}\fancyfoot[C]{\thepage}\fancyfoot[L]{\footnotesize\textcopyright\ Đặng Tấn Quí}\fancyfoot[R]{\footnotesize SĐT: 0377\,122\,210}\renewcommand{\headrulewidth}{0pt}\renewcommand{\footrulewidth}{0.4pt}}
\newtcolorbox{dinhnghia}[1][]{enhanced, breakable, colback=blue!3!white, colframe=blue!60!black, fonttitle=\bfseries, title={Định nghĩa}, #1}
\newtcolorbox{dinhly}[1][]{enhanced, breakable, colback=green!3!white, colframe=green!50!black, fonttitle=\bfseries, title={Định lý}, #1}
\newtcolorbox{tinhchat}[1][]{enhanced, breakable, colback=yellow!5!white, colframe=orange!50!black, fonttitle=\bfseries, title={Tính chất}, #1}
\newtcolorbox{phuongphap}[1][]{enhanced, breakable, colback=gray!5!white, colframe=gray!50!black, fonttitle=\bfseries, title={Phương pháp}, #1}
\newtcolorbox{luuy}[1][]{enhanced, breakable, colback=red!3!white, colframe=red!50!black, fonttitle=\bfseries, title={Lưu ý}, #1}
\newtcolorbox{congthuc}[1][]{enhanced, breakable, colback=cyan!3!white, colframe=cyan!50!black, fonttitle=\bfseries, title={Công thức}, #1}
\newtcolorbox{vidu}[1][]{enhanced, breakable, colback=violet!3!white, colframe=violet!50!black, fonttitle=\bfseries, title={Ví dụ}, #1}

\begin{document}
\thispagestyle{empty}
\begin{tikzpicture}[remember picture, overlay]
  \fill[blue!80!black] (current page.south west) rectangle (current page.north east);
  \node[anchor=west, white, font=\fontsize{26}{32}\bfseries\selectfont]
    at ([xshift=3cm, yshift=-7.5cm]current page.north west) {LÝ THUYẾT ĐỘ ĐO VÀ TÍCH PHÂN};
  \node[anchor=west, white, font=\fontsize{18}{24}\selectfont]
    at ([xshift=3cm, yshift=-9cm]current page.north west) {Lebesgue};
  \node[anchor=west, white, font=\Large\bfseries] at ([xshift=3cm, yshift=-13cm]current page.north west) {Biên soạn:};
  \node[anchor=west, yellow!80!white, font=\fontsize{22}{28}\bfseries\selectfont] at ([xshift=3cm, yshift=-14.5cm]current page.north west) {ĐẶNG TẤN QUÍ};
  \node[anchor=south east, white, font=\Large\bfseries] at ([xshift=-2cm, yshift=3cm]current page.south east) {\the\year};
\end{tikzpicture}
\newpage
\tableofcontents
\newpage
%% ================================================================
%%  CHƯƠNG 1: HỆ TẬP HỢP VÀ SIGMA-ĐẠI SỐ
%% ================================================================
\section{Hệ tập hợp và $\sigma$-đại số}

\subsection{Các phép toán trên tập hợp}

\begin{dinhnghia}[title={Giới hạn trên, giới hạn dưới của dãy tập hợp}]
  Cho dãy tập hợp $(A_n)_{n \ge 1}$:
  \begin{itemize}
    \item \textbf{Giới hạn trên:} $\displaystyle\limsup_{n \to \infty} A_n = \bigcap_{n=1}^{\infty} \bigcup_{k=n}^{\infty} A_k$ \;($x$ thuộc vô hạn tập $A_n$).
    \item \textbf{Giới hạn dưới:} $\displaystyle\liminf_{n \to \infty} A_n = \bigcup_{n=1}^{\infty} \bigcap_{k=n}^{\infty} A_k$ \;($x$ thuộc tất cả trừ hữu hạn tập $A_n$).
    \item Nếu $\limsup A_n = \liminf A_n = A$, ta nói dãy có \textbf{giới hạn} $A$, viết $\lim A_n = A$.
  \end{itemize}
\end{dinhnghia}

\subsection{$\sigma$-đại số}

\begin{dinhnghia}
  Cho $X \ne \varnothing$. Họ $\mathcal{F} \subset \mathcal{P}(X)$ là \textbf{$\sigma$-đại số} (sigma-algebra) trên $X$ nếu:
  \begin{enumerate}
    \item $X \in \mathcal{F}$.
    \item $A \in \mathcal{F} \;\Rightarrow\; A^c \in \mathcal{F}$ (đóng với phép lấy phần bù).
    \item $A_1, A_2, \ldots \in \mathcal{F} \;\Rightarrow\; \displaystyle\bigcup_{n=1}^{\infty} A_n \in \mathcal{F}$ (đóng với hợp đếm được).
  \end{enumerate}
  Cặp $(X, \mathcal{F})$ gọi là \textbf{không gian đo được} (measurable space).
\end{dinhnghia}

\begin{tinhchat}
  Nếu $\mathcal{F}$ là $\sigma$-đại số trên $X$ thì:
  \begin{itemize}
    \item $\varnothing \in \mathcal{F}$.
    \item $\mathcal{F}$ đóng với giao đếm được: $A_1, A_2, \ldots \in \mathcal{F} \Rightarrow \bigcap A_n \in \mathcal{F}$.
    \item $A, B \in \mathcal{F} \Rightarrow A \setminus B \in \mathcal{F}$.
    \item $\limsup A_n,\; \liminf A_n \in \mathcal{F}$.
  \end{itemize}
\end{tinhchat}

\begin{dinhnghia}[title={$\sigma$-đại số sinh bởi họ tập}]
  Cho $\mathcal{C} \subset \mathcal{P}(X)$. \textbf{$\sigma$-đại số sinh bởi $\mathcal{C}$}:
  \[
    \sigma(\mathcal{C}) = \bigcap\bigl\{\mathcal{F} : \mathcal{F} \text{ là } \sigma\text{-đại số trên } X,\; \mathcal{C} \subset \mathcal{F}\bigr\}.
  \]
  Đây là $\sigma$-đại số nhỏ nhất chứa $\mathcal{C}$.
\end{dinhnghia}

\begin{dinhnghia}[title={$\sigma$-đại số Borel}]
  Trên $\mathbb{R}^n$, \textbf{$\sigma$-đại số Borel} $\mathcal{B}(\mathbb{R}^n) = \sigma(\mathcal{O})$, trong đó $\mathcal{O}$ là họ tất cả các tập mở.

  Trên $\mathbb{R}$: $\mathcal{B}(\mathbb{R}) = \sigma\bigl(\{(a, b) : a < b\}\bigr) = \sigma\bigl(\{(-\infty, a] : a \in \mathbb{R}\}\bigr)$.
\end{dinhnghia}

\begin{vidu}
  \begin{itemize}
    \item $\sigma$-đại số tầm thường: $\{\varnothing, X\}$.
    \item $\sigma$-đại số lớn nhất: $\mathcal{P}(X)$.
    \item Tập Borel bao gồm: tập mở, tập đóng, $G_\delta$ (giao đếm tập mở), $F_\sigma$ (hợp đếm tập đóng), \ldots
  \end{itemize}
\end{vidu}

%% ================================================================
%%  CHƯƠNG 2: ĐỘ ĐO
%% ================================================================
\section{Độ đo}

\subsection{Định nghĩa và tính chất}

\begin{dinhnghia}
  Cho $(X, \mathcal{F})$ là không gian đo được. Hàm $\mu: \mathcal{F} \to [0, +\infty]$ là \textbf{độ đo} nếu:
  \begin{enumerate}
    \item $\mu(\varnothing) = 0$.
    \item \textbf{$\sigma$-cộng tính:} Nếu $A_1, A_2, \ldots \in \mathcal{F}$ đôi một rời nhau thì
      \[
        \mu\!\left(\bigcup_{n=1}^{\infty} A_n\right) = \sum_{n=1}^{\infty} \mu(A_n).
      \]
  \end{enumerate}
  Bộ ba $(X, \mathcal{F}, \mu)$ gọi là \textbf{không gian đo} (measure space).
\end{dinhnghia}

\begin{tinhchat}
  \begin{enumerate}
    \item \textbf{Đơn điệu:} $A \subset B \Rightarrow \mu(A) \le \mu(B)$.
    \item \textbf{Cộng dưới ($\sigma$-dưới cộng tính):} $\mu\!\left(\bigcup A_n\right) \le \sum \mu(A_n)$.
    \item \textbf{Hiệu:} Nếu $A \subset B$ và $\mu(A) < +\infty$ thì $\mu(B \setminus A) = \mu(B) - \mu(A)$.
    \item \textbf{Liên tục dưới:} $A_n \uparrow A$ (tăng) $\Rightarrow$ $\mu(A_n) \uparrow \mu(A)$.
    \item \textbf{Liên tục trên:} $A_n \downarrow A$ (giảm) và $\mu(A_1) < +\infty$ $\Rightarrow$ $\mu(A_n) \downarrow \mu(A)$.
  \end{enumerate}
\end{tinhchat}

\begin{dinhnghia}[title={Phân loại độ đo}]
  \begin{itemize}
    \item \textbf{Hữu hạn:} $\mu(X) < +\infty$.
    \item \textbf{$\sigma$-hữu hạn:} $X = \bigcup_{n=1}^{\infty} X_n$ với $\mu(X_n) < +\infty$.
    \item \textbf{Xác suất:} $\mu(X) = 1$.
    \item \textbf{Đầy đủ:} nếu $\mu(A) = 0$ và $B \subset A$ thì $B \in \mathcal{F}$ (và $\mu(B) = 0$).
  \end{itemize}
\end{dinhnghia}

\subsection{Độ đo ngoài và sự mở rộng}

\begin{dinhnghia}[title={Độ đo ngoài}]
  $\mu^*: \mathcal{P}(X) \to [0, +\infty]$ là \textbf{độ đo ngoài} nếu:
  \begin{enumerate}
    \item $\mu^*(\varnothing) = 0$.
    \item Đơn điệu: $A \subset B \Rightarrow \mu^*(A) \le \mu^*(B)$.
    \item $\sigma$-dưới cộng tính: $\mu^*\!\left(\bigcup A_n\right) \le \sum \mu^*(A_n)$.
  \end{enumerate}
\end{dinhnghia}

\begin{dinhly}[title={Định lý Carathéodory}]
  Cho $\mu^*$ là độ đo ngoài trên $X$. Đặt:
  \[
    \mathcal{M} = \{A \subset X : \mu^*(E) = \mu^*(E \cap A) + \mu^*(E \cap A^c),\; \forall E \subset X\}.
  \]
  Thì $\mathcal{M}$ là $\sigma$-đại số và $\mu^*|_{\mathcal{M}}$ là độ đo đầy đủ.
\end{dinhly}

\begin{dinhly}[title={Định lý mở rộng Carathéodory--Hahn}]
  Cho $\mu_0$ là tiền độ đo trên đại số $\mathcal{A}$. Thì tồn tại độ đo $\mu$ trên $\sigma(\mathcal{A})$ sao cho $\mu|_{\mathcal{A}} = \mu_0$. Nếu $\mu_0$ là $\sigma$-hữu hạn thì sự mở rộng là duy nhất.
\end{dinhly}

\subsection{Độ đo Lebesgue}

\begin{dinhnghia}[title={Độ đo Lebesgue ngoài trên $\mathbb{R}$}]
  Với $A \subset \mathbb{R}$:
  \[
    m^*(A) = \inf\left\{\sum_{n=1}^{\infty} (b_n - a_n) : A \subset \bigcup_{n=1}^{\infty} (a_n, b_n)\right\}.
  \]
\end{dinhnghia}

\begin{dinhnghia}[title={Tập đo được Lebesgue}]
  $A \subset \mathbb{R}$ \textbf{đo được Lebesgue} nếu $A$ thỏa điều kiện Carathéodory đối với $m^*$. Họ các tập đo được Lebesgue ký hiệu $\mathcal{L}$, và $m = m^*|_{\mathcal{L}}$ là \textbf{độ đo Lebesgue}.
\end{dinhnghia}

\begin{tinhchat}[title={Tính chất độ đo Lebesgue}]
  \begin{enumerate}
    \item $\mathcal{B}(\mathbb{R}) \subset \mathcal{L}$ (mọi tập Borel đều đo được Lebesgue).
    \item $m\bigl((a, b)\bigr) = m\bigl([a, b]\bigr) = b - a$.
    \item Tập đếm được có độ đo $0$.
    \item \textbf{Bất biến tịnh tiến:} $m(A + x) = m(A)$.
    \item \textbf{Co giãn:} $m(cA) = |c|\,m(A)$.
    \item $(\mathbb{R}, \mathcal{L}, m)$ là không gian đo $\sigma$-hữu hạn và đầy đủ.
    \item Tồn tại tập không đo được Lebesgue (ví dụ: tập Vitali).
  \end{enumerate}
\end{tinhchat}

\begin{dinhnghia}[title={Độ đo Lebesgue trên $\mathbb{R}^n$}]
  Với hộp $I = (a_1, b_1) \times \cdots \times (a_n, b_n)$: $\text{vol}(I) = \prod_{i=1}^{n}(b_i - a_i)$.

  Độ đo Lebesgue ngoài: $m^*(A) = \inf\bigl\{\sum \text{vol}(I_k) : A \subset \bigcup I_k\bigr\}$.

  Định lý Carathéodory cho $\sigma$-đại số và độ đo Lebesgue $m$ trên $\mathbb{R}^n$.
\end{dinhnghia}

\begin{dinhly}[title={Xấp xỉ tập đo được}]
  Với $A \in \mathcal{L}$, $m(A) < +\infty$, $\forall \varepsilon > 0$:
  \begin{enumerate}
    \item Tồn tại tập mở $G \supset A$: $m(G \setminus A) < \varepsilon$.
    \item Tồn tại tập đóng $F \subset A$: $m(A \setminus F) < \varepsilon$.
    \item Tồn tại hợp hữu hạn khoảng $U$: $m(A \triangle U) < \varepsilon$.
  \end{enumerate}
\end{dinhly}

%% ================================================================
%%  CHƯƠNG 3: HÀM ĐO ĐƯỢC
%% ================================================================
\section{Hàm đo được}

\begin{dinhnghia}
  Cho $(X, \mathcal{F})$ và $(\mathbb{R}, \mathcal{B})$. Hàm $f: X \to \overline{\mathbb{R}} = \mathbb{R} \cup \{-\infty, +\infty\}$ là \textbf{đo được} (đối với $\mathcal{F}$) nếu:
  \[
    \forall a \in \mathbb{R}:\; \{x \in X : f(x) > a\} \in \mathcal{F}.
  \]
\end{dinhnghia}

\begin{tinhchat}
  Các điều kiện tương đương ($f$ đo được):
  \begin{itemize}
    \item $\{f > a\} \in \mathcal{F}$, $\forall a$.
    \item $\{f \ge a\} \in \mathcal{F}$, $\forall a$.
    \item $\{f < a\} \in \mathcal{F}$, $\forall a$.
    \item $\{f \le a\} \in \mathcal{F}$, $\forall a$.
    \item $f^{-1}(B) \in \mathcal{F}$, $\forall B \in \mathcal{B}(\mathbb{R})$.
  \end{itemize}
\end{tinhchat}

\begin{tinhchat}[title={Phép toán trên hàm đo được}]
  Nếu $f, g$ đo được, $c \in \mathbb{R}$ thì:
  \begin{itemize}
    \item $cf$, \;$f + g$, \;$fg$, \;$f/g$ (khi $g \ne 0$) đo được.
    \item $|f|$, \;$f^+ = \max(f, 0)$, \;$f^- = \max(-f, 0)$ đo được; $f = f^+ - f^-$.
    \item $\max(f, g)$, \;$\min(f, g)$ đo được.
    \item $\sup_n f_n$, \;$\inf_n f_n$, \;$\limsup f_n$, \;$\liminf f_n$ đo được.
  \end{itemize}
\end{tinhchat}

\begin{dinhly}
  Nếu $(f_n)$ là dãy hàm đo được và $f_n \to f$ điểm (hay hầu khắp nơi) thì $f$ đo được.
\end{dinhly}

\begin{dinhnghia}[title={Hàm đơn giản (simple function)}]
  $\varphi: X \to \mathbb{R}$ là \textbf{hàm đơn giản} nếu $\varphi$ chỉ nhận hữu hạn giá trị:
  \[
    \varphi = \sum_{i=1}^{n} c_i\, \chi_{A_i},
  \]
  trong đó $A_i = \{x : \varphi(x) = c_i\} \in \mathcal{F}$ và $\chi_A$ là hàm đặc trưng của $A$.
\end{dinhnghia}

\begin{dinhly}[title={Xấp xỉ bởi hàm đơn giản}]
  Nếu $f \ge 0$ đo được thì tồn tại dãy hàm đơn giản $(\varphi_n)$ tăng sao cho $\varphi_n \uparrow f$ điểm. Nếu $f$ bất kỳ (đo được) thì tồn tại $(\varphi_n)$ đo được, $\varphi_n \to f$ điểm và $|\varphi_n| \le |f|$.
\end{dinhly}

\begin{dinhnghia}[title={Hầu khắp nơi (almost everywhere --- a.e.)}]
  Tính chất $P(x)$ đúng \textbf{hầu khắp nơi} (viết tắt h.k.n. hay a.e.) nếu $\mu\bigl(\{x : P(x) \text{ sai}\}\bigr) = 0$.
\end{dinhnghia}

\begin{dinhly}[title={Định lý Egorov}]
  Cho $\mu(X) < +\infty$ và $f_n \to f$ a.e. Thì $\forall \varepsilon > 0$, tồn tại $E \in \mathcal{F}$ với $\mu(E) < \varepsilon$ sao cho $f_n \to f$ \textbf{đều} trên $X \setminus E$.
\end{dinhly}

\begin{dinhly}[title={Định lý Lusin}]
  Cho $f: \mathbb{R} \to \mathbb{R}$ đo được Lebesgue, $m(A) < +\infty$. Thì $\forall \varepsilon > 0$, tồn tại tập đóng $F \subset A$ với $m(A \setminus F) < \varepsilon$ sao cho $f|_F$ liên tục.
\end{dinhly}

%% ================================================================
%%  CHƯƠNG 4: TÍCH PHÂN LEBESGUE
%% ================================================================
\section{Tích phân Lebesgue}

\subsection{Định nghĩa tích phân}

\begin{dinhnghia}[title={Tích phân hàm đơn giản}]
  Với $\varphi = \sum_{i=1}^{n} c_i\, \chi_{A_i} \ge 0$:
  \[
    \int_X \varphi\,d\mu = \sum_{i=1}^{n} c_i\, \mu(A_i).
  \]
\end{dinhnghia}

\begin{dinhnghia}[title={Tích phân hàm không âm}]
  Với $f \ge 0$ đo được:
  \[
    \int_X f\,d\mu = \sup\left\{\int_X \varphi\,d\mu : 0 \le \varphi \le f,\; \varphi \text{ đơn giản}\right\}.
  \]
\end{dinhnghia}

\begin{dinhnghia}[title={Tích phân hàm tổng quát}]
  $f$ đo được, viết $f = f^+ - f^-$. $f$ \textbf{khả tích Lebesgue} (hay $f \in L^1$) nếu $\displaystyle\int f^+\,d\mu < +\infty$ và $\displaystyle\int f^-\,d\mu < +\infty$. Khi đó:
  \[
    \int_X f\,d\mu = \int_X f^+\,d\mu - \int_X f^-\,d\mu.
  \]
  Tương đương: $f$ khả tích $\Leftrightarrow$ $\displaystyle\int |f|\,d\mu < +\infty$.
\end{dinhnghia}

\begin{luuy}
  Ký hiệu: $\displaystyle\int_A f\,d\mu = \int_X f\,\chi_A\,d\mu$. Khi $\mu$ là độ đo Lebesgue trên $\mathbb{R}$: $\displaystyle\int_a^b f\,dm = \int_{[a,b]} f\,dm$.
\end{luuy}

\subsection{Tính chất tích phân Lebesgue}

\begin{tinhchat}
  Cho $f, g$ khả tích:
  \begin{enumerate}
    \item \textbf{Tuyến tính:} $\displaystyle\int (\alpha f + \beta g)\,d\mu = \alpha\int f\,d\mu + \beta\int g\,d\mu$.
    \item \textbf{Đơn điệu:} $f \le g$ a.e. $\Rightarrow$ $\displaystyle\int f\,d\mu \le \int g\,d\mu$.
    \item $\left|\displaystyle\int f\,d\mu\right| \le \int |f|\,d\mu$.
    \item $f = g$ a.e. $\Rightarrow$ $\displaystyle\int f\,d\mu = \int g\,d\mu$.
    \item $f \ge 0$ và $\displaystyle\int f\,d\mu = 0$ $\Rightarrow$ $f = 0$ a.e.
    \item \textbf{$\sigma$-cộng tính trên tập:} Nếu $A = \bigsqcup A_n$ thì $\displaystyle\int_A f\,d\mu = \sum_n \int_{A_n} f\,d\mu$.
    \item $f$ khả tích $\Rightarrow$ $\{|f| = +\infty\}$ có độ đo $0$.
  \end{enumerate}
\end{tinhchat}

\begin{dinhly}[title={So sánh Riemann và Lebesgue}]
  Nếu $f$ khả tích Riemann trên $[a, b]$ thì $f$ khả tích Lebesgue trên $[a, b]$ và hai tích phân bằng nhau:
  \[
    (R)\!\int_a^b f(x)\,dx = \int_{[a,b]} f\,dm.
  \]
  Điều ngược lại không đúng: $f$ có thể khả tích Lebesgue nhưng không khả tích Riemann.
\end{dinhly}

\subsection{Các định lý hội tụ}

\begin{dinhly}[title={Bổ đề Fatou}]
  Cho $(f_n)$ là dãy hàm đo được, $f_n \ge 0$ a.e. Thì:
  \[
    \int \liminf_{n \to \infty} f_n\,d\mu \le \liminf_{n \to \infty} \int f_n\,d\mu.
  \]
\end{dinhly}

\begin{dinhly}[title={Định lý hội tụ đơn điệu (MCT --- Beppo Levi)}]
  Cho $(f_n)$ đo được, $0 \le f_1 \le f_2 \le \cdots$ a.e. và $f_n \uparrow f$ a.e. Thì:
  \[
    \lim_{n \to \infty} \int f_n\,d\mu = \int f\,d\mu \quad(\text{có thể } = +\infty).
  \]
\end{dinhly}

\begin{dinhly}[title={Định lý hội tụ bị chặn (DCT --- Lebesgue)}]
  Cho $(f_n)$ đo được, $f_n \to f$ a.e. Nếu tồn tại $g \ge 0$ \textbf{khả tích} sao cho $|f_n| \le g$ a.e. $\forall n$, thì:
  \[
    \lim_{n \to \infty} \int f_n\,d\mu = \int f\,d\mu, \qquad\text{và}\qquad \lim_{n \to \infty}\int |f_n - f|\,d\mu = 0.
  \]
\end{dinhly}

\begin{luuy}
  DCT là công cụ mạnh nhất để đổi giới hạn và tích phân. Điều kiện ``bị chặn bởi hàm khả tích'' là thiết yếu.
\end{luuy}

\begin{dinhly}[title={Hệ quả: chuỗi hàm}]
  Nếu $f_n \ge 0$ đo được thì:
  \[
    \int \sum_{n=1}^{\infty} f_n\,d\mu = \sum_{n=1}^{\infty} \int f_n\,d\mu.
  \]
\end{dinhly}

\begin{dinhly}[title={Đạo hàm dưới dấu tích phân}]
  Cho $f(x, t)$ đo được theo $x$, khả vi theo $t$. Nếu $|f'_t(x, t)| \le g(x)$ với $g$ khả tích, thì:
  \[
    \frac{d}{dt}\int f(x, t)\,d\mu(x) = \int \frac{\partial f}{\partial t}(x, t)\,d\mu(x).
  \]
\end{dinhly}

%% ================================================================
%%  CHƯƠNG 5: KHÔNG GIAN $L^p$
%% ================================================================
\section{Không gian $L^p$}

\begin{dinhnghia}
  Cho $(X, \mathcal{F}, \mu)$ và $1 \le p < +\infty$:
  \[
    L^p(X, \mu) = \left\{f \text{ đo được} : \int_X |f|^p\,d\mu < +\infty\right\}.
  \]
  \textbf{Chuẩn $L^p$:} $\|f\|_p = \left(\displaystyle\int |f|^p\,d\mu\right)^{1/p}$.

  $L^\infty(X, \mu) = \{f \text{ đo được} : \|f\|_\infty < +\infty\}$, với $\|f\|_\infty = \text{ess\,sup}\,|f| = \inf\{M : |f| \le M \text{ a.e.}\}$.

  Trong $L^p$, ta đồng nhất $f \sim g$ nếu $f = g$ a.e.
\end{dinhnghia}

\begin{dinhly}[title={Bất đẳng thức Hölder}]
  Cho $1 \le p, q \le +\infty$ với $\dfrac{1}{p} + \dfrac{1}{q} = 1$. Nếu $f \in L^p$ và $g \in L^q$ thì $fg \in L^1$ và:
  \[
    \int |fg|\,d\mu \le \|f\|_p\,\|g\|_q.
  \]
  Trường hợp $p = q = 2$: bất đẳng thức Cauchy--Schwarz.
\end{dinhly}

\begin{dinhly}[title={Bất đẳng thức Minkowski}]
  Với $1 \le p \le +\infty$, nếu $f, g \in L^p$ thì $f + g \in L^p$ và:
  \[
    \|f + g\|_p \le \|f\|_p + \|g\|_p.
  \]
\end{dinhly}

\begin{dinhly}[title={Tính đầy đủ (Riesz--Fischer)}]
  $L^p(X, \mu)$ là \textbf{không gian Banach} (không gian định chuẩn đầy đủ) với $1 \le p \le +\infty$.

  $L^2(X, \mu)$ là \textbf{không gian Hilbert} với tích vô hướng $\langle f, g \rangle = \displaystyle\int f\,\overline{g}\,d\mu$.
\end{dinhly}

\begin{dinhly}[title={Hội tụ trong $L^p$}]
  $f_n \to f$ trong $L^p$ (tức $\|f_n - f\|_p \to 0$) thì:
  \begin{itemize}
    \item Tồn tại dãy con $f_{n_k} \to f$ a.e.
    \item $\|f_n\|_p \to \|f\|_p$.
  \end{itemize}
  Chú ý: hội tụ trong $L^p$ không suy ra hội tụ a.e., và ngược lại.
\end{dinhly}

\begin{dinhly}[title={Tập trù mật}]
  \begin{itemize}
    \item Hàm đơn giản trù mật trong $L^p$ ($1 \le p < \infty$).
    \item $C_c(\mathbb{R}^n)$ (hàm liên tục có giá compact) trù mật trong $L^p(\mathbb{R}^n, m)$.
    \item $L^p(\mathbb{R}^n)$ là \textbf{khả li} (separable) khi $1 \le p < \infty$.
  \end{itemize}
\end{dinhly}

%% ================================================================
%%  CHƯƠNG 6: ĐỘ ĐO TÍCH VÀ ĐỊNH LÝ FUBINI
%% ================================================================
\section{Độ đo tích và định lý Fubini}

\begin{dinhnghia}[title={Độ đo tích}]
  Cho $(X, \mathcal{F}, \mu)$ và $(Y, \mathcal{G}, \nu)$. \textbf{$\sigma$-đại số tích}: $\mathcal{F} \otimes \mathcal{G} = \sigma\{A \times B : A \in \mathcal{F},\, B \in \mathcal{G}\}$.

  \textbf{Độ đo tích} $\mu \otimes \nu$ trên $\mathcal{F} \otimes \mathcal{G}$: xác định duy nhất bởi $(\mu \otimes \nu)(A \times B) = \mu(A)\,\nu(B)$.
\end{dinhnghia}

\begin{dinhly}[title={Định lý Tonelli (hàm không âm)}]
  Cho $f: X \times Y \to [0, +\infty]$ đo được đối với $\mathcal{F} \otimes \mathcal{G}$. Thì:
  \[
    \int_{X \times Y} f\,d(\mu \otimes \nu) = \int_X\!\left(\int_Y f(x,y)\,d\nu(y)\right)d\mu(x) = \int_Y\!\left(\int_X f(x,y)\,d\mu(x)\right)d\nu(y).
  \]
\end{dinhly}

\begin{dinhly}[title={Định lý Fubini}]
  Nếu $f \in L^1(X \times Y, \mu \otimes \nu)$ thì:
  \begin{enumerate}
    \item $f(x, \cdot) \in L^1(Y, \nu)$ cho $\mu$-hầu hết $x$; \;$f(\cdot, y) \in L^1(X, \mu)$ cho $\nu$-hầu hết $y$.
    \item $g(x) = \displaystyle\int_Y f(x,y)\,d\nu(y)$ khả tích theo $\mu$; \;$h(y) = \displaystyle\int_X f(x,y)\,d\mu(x)$ khả tích theo $\nu$.
    \item $\displaystyle\int_{X \times Y} f\,d(\mu \otimes \nu) = \int_X g\,d\mu = \int_Y h\,d\nu$.
  \end{enumerate}
\end{dinhly}

\begin{luuy}
  Thực hành: dùng Tonelli trước để kiểm tra khả tích (vì Tonelli áp dụng cho $|f| \ge 0$ không cần điều kiện khả tích), rồi áp dụng Fubini.
\end{luuy}

%% ================================================================
%%  CHƯƠNG 7: ĐỘ ĐO CÓ DẤU VÀ ĐỊNH LÝ RADON--NIKODYM
%% ================================================================
\section{Độ đo có dấu và định lý Radon--Nikodym}

\subsection{Độ đo có dấu}

\begin{dinhnghia}
  $\nu: \mathcal{F} \to [-\infty, +\infty]$ là \textbf{độ đo có dấu} (signed measure) nếu:
  \begin{enumerate}
    \item $\nu(\varnothing) = 0$.
    \item $\nu$ nhận nhiều nhất một trong hai giá trị $+\infty$, $-\infty$.
    \item $\sigma$-cộng tính: $\nu\bigl(\bigsqcup A_n\bigr) = \sum \nu(A_n)$ (chuỗi hội tụ tuyệt đối khi vế trái hữu hạn).
  \end{enumerate}
\end{dinhnghia}

\begin{dinhly}[title={Phân tích Hahn}]
  Cho $\nu$ là độ đo có dấu trên $(X, \mathcal{F})$. Thì tồn tại $P, N \in \mathcal{F}$ sao cho:
  \begin{itemize}
    \item $X = P \sqcup N$ (phân hoạch).
    \item $P$ là tập dương: $\nu(A) \ge 0$ với mọi $A \subset P$ đo được.
    \item $N$ là tập âm: $\nu(A) \le 0$ với mọi $A \subset N$ đo được.
  \end{itemize}
\end{dinhly}

\begin{dinhly}[title={Phân tích Jordan}]
  Mọi độ đo có dấu $\nu$ phân tích duy nhất thành $\nu = \nu^+ - \nu^-$, trong đó $\nu^+, \nu^-$ là độ đo (không âm) và kỳ dị lẫn nhau: $\nu^+(A) = \nu(A \cap P)$, $\nu^-(A) = -\nu(A \cap N)$.

  \textbf{Biến phân toàn phần:} $|\nu| = \nu^+ + \nu^-$.
\end{dinhly}

\subsection{Liên tục tuyệt đối và kỳ dị}

\begin{dinhnghia}
  Cho độ đo $\mu$ và $\nu$ trên $(X, \mathcal{F})$:
  \begin{itemize}
    \item $\nu$ \textbf{liên tục tuyệt đối} đối với $\mu$ ($\nu \ll \mu$) nếu: $\mu(A) = 0 \Rightarrow \nu(A) = 0$.
    \item $\nu$ \textbf{kỳ dị} đối với $\mu$ ($\nu \perp \mu$) nếu tồn tại $A \in \mathcal{F}$: $\mu(A) = 0$ và $\nu(A^c) = 0$.
  \end{itemize}
\end{dinhnghia}

\begin{dinhly}[title={Định lý Radon--Nikodym}]
  Cho $\mu$ là độ đo $\sigma$-hữu hạn và $\nu \ll \mu$. Thì tồn tại duy nhất (a.e.) hàm đo được $f \ge 0$ sao cho:
  \[
    \nu(A) = \int_A f\,d\mu, \qquad \forall A \in \mathcal{F}.
  \]
  Hàm $f$ gọi là \textbf{đạo hàm Radon--Nikodym}, ký hiệu $f = \dfrac{d\nu}{d\mu}$.
\end{dinhly}

\begin{dinhly}[title={Phân tích Lebesgue}]
  Cho $\mu$ $\sigma$-hữu hạn và $\nu$ $\sigma$-hữu hạn. Thì tồn tại duy nhất:
  \[
    \nu = \nu_a + \nu_s,
  \]
  trong đó $\nu_a \ll \mu$ và $\nu_s \perp \mu$.
\end{dinhly}

\begin{congthuc}[title={Quy tắc dây chuyền}]
  Nếu $\nu \ll \mu \ll \lambda$ thì $\nu \ll \lambda$ và:
  \[
    \frac{d\nu}{d\lambda} = \frac{d\nu}{d\mu} \cdot \frac{d\mu}{d\lambda} \quad\text{a.e. } [\lambda].
  \]
\end{congthuc}

%% ================================================================
%%  CHƯƠNG 8: ĐỘ ĐO TRÊN KHÔNG GIAN METRIC --- BỔ SUNG
%% ================================================================
\section{Bổ sung}

\subsection{Các loại hội tụ}

\begin{dinhnghia}
  Cho $(f_n)$ dãy hàm đo được trên $(X, \mathcal{F}, \mu)$:
  \begin{itemize}
    \item \textbf{Hội tụ điểm}: $f_n(x) \to f(x)$ với mọi $x$.
    \item \textbf{Hội tụ hầu khắp nơi (a.e.)}: $f_n \to f$ a.e., tức $\mu(\{x : f_n(x) \not\to f(x)\}) = 0$.
    \item \textbf{Hội tụ theo độ đo}: $\forall \varepsilon > 0$: $\mu(\{|f_n - f| > \varepsilon\}) \to 0$.
    \item \textbf{Hội tụ trong $L^p$}: $\|f_n - f\|_p \to 0$.
    \item \textbf{Hội tụ hầu đều}: $\forall \varepsilon > 0$, $\exists E$ với $\mu(E) < \varepsilon$ sao cho $f_n \to f$ đều trên $X \setminus E$.
  \end{itemize}
\end{dinhnghia}

\begin{dinhly}[title={Quan hệ giữa các loại hội tụ}]
  \begin{itemize}
    \item Hội tụ a.e. + $\mu(X) < \infty$ $\Rightarrow$ hội tụ theo độ đo (nhưng không ngược lại).
    \item Hội tụ trong $L^p$ $\Rightarrow$ hội tụ theo độ đo.
    \item Hội tụ theo độ đo $\Rightarrow$ tồn tại dãy con hội tụ a.e.
    \item Hội tụ a.e. + $\mu(X) < \infty$ $\Leftrightarrow$ hội tụ hầu đều (Egorov).
  \end{itemize}
\end{dinhly}

\subsection{Tích chập}

\begin{dinhnghia}
  Với $f, g \in L^1(\mathbb{R}^n)$, \textbf{tích chập}:
  \[
    (f * g)(x) = \int_{\mathbb{R}^n} f(x - y)\,g(y)\,dy.
  \]
\end{dinhnghia}

\begin{tinhchat}
  \begin{itemize}
    \item $f * g = g * f$ (giao hoán).
    \item $\|f * g\|_1 \le \|f\|_1 \cdot \|g\|_1$ (Young, trường hợp $p = q = 1$).
    \item $f * g$ liên tục nếu $f \in L^1$ và $g \in L^\infty$ (hoặc ngược lại).
  \end{itemize}
\end{tinhchat}

\begin{dinhly}[title={Bất đẳng thức Young (tổng quát)}]
  Nếu $\dfrac{1}{p} + \dfrac{1}{q} = 1 + \dfrac{1}{r}$ ($1 \le p, q, r \le \infty$), $f \in L^p$, $g \in L^q$ thì:
  \[
    \|f * g\|_r \le \|f\|_p \cdot \|g\|_q.
  \]
\end{dinhly}

\end{document}